% from ega4-15-fr.tex
\oldpage[IV-3]{226}
En effet, une composante irréductible $X'$ de $X$ contenant $x$ et telle que
$\dim_x(X'\cap f^{-1}(y))=\dim_x(f^{-1}(y))$ contient nécessairement $Z$,
puisqu'une composante irréductible de $X'\cap f^{-1}(y)$ contenant $x$ doit être
de même dimension que la seule composante irréductible $Z$ de $f^{-1}(y)$
contenant $x$. On peut alors appliquer (15.1.2) à $z$.

\begin{remarks}[15.1.4]
\label{IV.15.1.4-fr}
\begin{enumerate}[label=(\roman*)]
  \item Dans l'énoncé de (15.1.1), on ne peut remplacer «~universellement
  ouvert~» par «~équidimensionnel~», comme le montre l'exemple (14.4.10, (ii)) où
  l'on fait $\mathcal F=\mathcal O_X$~; les fibres de $f$ sont alors des schémas
  artiniens réduits, donc (avec les notations introduites \emph{loc. cit.}) on a
  $\lambda_{x'}(\mathcal F_y)=1$, mais il y a deux composantes irréductibles
  $X_1$, $X_2$ de $X$ contenant $x'$, et le second membre de (15.1.1.1) est donc
  égal à $2$, bien que la restriction de $f$ à chacune des composantes $X_1$,
  $X_2$ soit un morphisme équidimensionnel en tout point. Ce même exemple montre
  aussi que dans (15.1.2) et (15.1.3), on ne peut remplacer l'hypothèse
  «~universellement ouvert~» par «~équidimensionnel~».
  \item Il est vraisemblable que pour la validité de l'inégalité (15.1.1.2), on
  ne peut supprimer l'hypothèse que $\mathcal O_y$ est \emph{régulier}.
\end{enumerate}
\end{remarks}

\subsection{Platitude des morphismes universellement ouverts à fibres géométriquement réduites}
\label{subsection:IV.15.2-fr}

\begin{env}[15.2.1]
\label{IV.15.2.1-fr}
On a vu (2.4.6) qu'un morphisme \emph{plat} et localement de présentation finie
est \emph{universellement ouvert}. On a une réciproque partielle de ce résultat
moyennant des hypothèses supplémentaires~:
\end{env}

\begin{theorem}[15.2.2]
\label{IV.15.2.2-fr}
Soient $Y$ un préschéma localement noethérien, $f:X\to Y$ un morphisme localement
de type fini, $\mathcal F$ un $\mathcal O_X$-Module cohérent, $x$ un point de
$X$, $y=f(x)$. On suppose vérifiées les conditions suivantes~:
\begin{enumerate}[label=(\roman*)]
  \item $\operatorname{Supp}(\mathcal F)=X$.
  \item $f$ est universellement ouvert aux points génériques des composantes
  irréductibles de $f^{-1}(y)$ contenant $x$.
  \item $\mathcal F_y=\mathcal F\otimes_{\mathcal O_Y}\mathbf{k}(y)$ est
  géométriquement réduit au point $x$ (4.6.22).
  \item L'anneau $\mathcal O_y$ est réduit.
\end{enumerate}
Alors $\mathcal F$ est $f$-plat au point $x$.
\end{theorem}

\begin{proof}
Comme $\mathcal O_y$ est réduit et noethérien, nous allons appliquer le critère
valuatif de platitude (11.8.1). Considérons donc un homomorphisme local
$\mathcal O_y\to A'$, où $A'$ est un anneau de valuation discrète~; posons
$Y'=\operatorname{Spec}(A')$, $X'=X\times_Y Y'$, $f'=f_{(Y')}:X'\to Y'$, et
soit $x'$ un point de $X'$ dont les projections dans $X$ et $Y'$ sont
respectivement $x$ et le point fermé $y'$ de $Y'$~; si l'on pose
$\mathcal F'=\mathcal F\otimes_Y Y'$, on a
$\operatorname{Supp}(\mathcal F')=X'$ par (\textbf{I}, 9.1.13)~; tout point
générique $z'$ d'une composante irréductible de $f'^{-1}(y')$ contenant $x'$ a
pour projection dans $f^{-1}(y)$ un point générique $z$ d'une composante
irréductible de $f^{-1}(y)$ contenant $x$, en vertu de (4.2.6)~; donc $f'$ est
universellement ouvert au point $z'$. Enfin, il résulte de (4.7.11) que
$\mathcal F'_{y'}$ est géométriquement réduit au point $x'$. On voit donc qu'on
est ramené à démontrer le

